\section{Conclusion and Future Work}
\label{sec:conclusion}

This paper has presented a Compute-Continuum-aware workflow orchestration framework that integrates data- and bandwidth-aware scheduling, orchestration-level decision policy, and FAIR-compatible provenance for eScience workflows. Its optimisation core---a mixed-integer formulation with correctness proofs, together with a randomized relaxation, a greedy scheduler and a genetic metaheuristic---is carried over from our conference study~\cite{gupta2025efficient}; the framework built around that core, and the positioning of federated fog as the latency-sensitive tier of a broader execution continuum, are contributions of this submission. It interprets the federated fog as the latency-sensitive tier of a broader execution environment that can include edge, cloud, and HPC resources when their capabilities and communication paths are compatible with workflow requirements.

The reported comparisons show how the scheduling alternatives respond to four dimensions relevant to Compute-Continuum orchestration: deadline flexibility, federated resource scale, workflow scale, and workflow-node failures. The exact MILP serves as a global optimisation benchmark, whereas Relaxation and the Metaheuristic provide scalable alternatives. The Greedy method offers a lower-overhead baseline that remains sensitive to task dependencies and communication delay. The result is a clear decision framework: use exact optimisation when the instance is small enough, use relaxation or metaheuristics when workflow scale or resource volatility makes exact solution impractical, and use the greedy policy when a fast local decision is required.

The paper also identifies reproducibility as a first-class workflow requirement. A reusable scheduling artifact should expose the workflow DAG, data-size assumptions, resource graph, configuration profile, randomisation settings, and run-level provenance. This information enables another investigator to reproduce a placement decision and trace a reported figure to its inputs. The proposed FAIR-oriented methodology is intentionally implementation-neutral so it can be adopted by different workflow-management and orchestration systems.

Several limitations define the next research steps. The supplied results are simulation-based and do not establish measured energy consumption, cloud/HPC execution cost, security compliance, or performance on a particular exascale platform. Future work should therefore deploy the artifact across real edge, fog, cloud, and HPC resources; collect trace-level measurements; evaluate energy- and cost-aware policies; and test recovery mechanisms under network partition and node failure. Further extensions can investigate dynamic workload adaptation, graph compression, hybrid partitioning strategies, and uncertainty-aware accuracy propagation through scientific workflows. Taken together, the formulation, the orchestration policy, and the reproducibility methodology define a scheduling contribution that a Compute-Continuum workflow system can adopt directly, and a set of concrete deployment steps---artifact deposition, standard-format serialisation, and trace-level validation across real edge, fog, cloud and HPC resources---that would carry it from a validated model to a validated system.
